\input{../tex-inputs/latex-leseansicht-vorspann}

\section[Theodor von Sosnosky an Arthur Schnitzler, 23. 5. 1901]{L04230 Theodor von Sosnosky an Arthur Schnitzler, 23. 5. 1901}
\nopagebreak\mylabel{L04230v}
\rehead{ }\normalsize\beginnumbering\briefempfaengerindex{Schnitzler, Arthur@\textsc{Schnitzler, Arthur}!zzzSosnosky, Theodor von@\emph{von Theodor von Sosnosky}!1901-05-231@{23. 5. 1901}|(be}
\toendnotes[C]{\smallbreak\pagebreak[2]}
\correspDesc{Versand  durch Theodor von Sosnosky am 23. 5. 1901 in Kremsmünster
\newline{}Erhalt  durch Arthur Schnitzler im Zeitraum [24. 5. 1901
                  – 28. 5. 1901?] in Wien}\toendnotes[C]{\smallbreak}
\Standort{DLA, A:Schnitzler, HS.1985.1.4640.}
\physDesc{Briefkarte, 781 Zeichen
\newline{}Handschrift: blaue Tinte, deutsche Kurrent}\toendnotes[C]{\smallbreak}
\pstart
           {\pb}\label{T_L04230-1v}\edtext{dzt \textsc{Kr}\oindex{Kremsmünster@\textbf{Kremsmünster}, \emph{Verwaltungsgebiet}|pw}{ }23/5 901}{\lemma{\textnormal{\emph{dzt Kr 23/5 901}}}\Cendnote{\textnormal{Die Datumszeile befindet sich am
                     unteren Rand von S. 1.}}}\label{T_L04230-1}\pend
           \vspace{0.5em}
\pstart
           Sehr geehrter Herr Doktor! Vertraulichſten Dank für die
               liebenswürdige Zuſendung Ihres \label{K_L04230-1v}\edtext{jüngſten Werkes\pwindex{Schnitzler, Arthur 15.\,5.\,1862 Wien – 21.\,10.\,1931 ebd.@\textsc{Schnitzler, Arthur} (15.\,5.\,1862 Wien – 21.\,10.\,1931 ebd.), \emph{Schriftsteller, Mediziner}!Lieutenant Gustl. Novelle@\strich\emph{Lieutenant Gustl. Novelle}|pwv}}{\lemma{\textnormal{\emph{jüngsten Werkes}}}\Cendnote{\textnormal{Die Buchausgabe von \emph{Lieutnant Gustl}\pwindex{Schnitzler, Arthur 15.\,5.\,1862 Wien – 21.\,10.\,1931 ebd.@\textsc{Schnitzler, Arthur} (15.\,5.\,1862 Wien – 21.\,10.\,1931 ebd.), \emph{Schriftsteller, Mediziner}!Lieutenant Gustl. Novelle@\strich\emph{Lieutenant Gustl. Novelle}|pwk} wurde am 23. 5. 1901 vom \emph{Börsenblatt für den deutschen Buchhandel}\pwindex{Börsenblatt für den Deutschen Buchhandel@\emph{Börsenblatt für den Deutschen Buchhandel}|pwk} als
                  Neuerscheinung gemeldet.}}}\label{K_L04230-1}! Gerade geſtern hab’ ich die Beſprechung\pwindex{Sosnosky, Theodor von 4.\,1.\,1866 Budapest – 3.\,2.\,1943 Wien@\textsc{Sosnosky, Theodor von} (4.\,1.\,1866 Budapest – 3.\,2.\,1943 Wien), \emph{Schriftsteller}!Bücher und Zeitschriftenschau [Frau Bertha Garlan]@\strich\emph{Bücher und Zeitschriftenschau [Frau Bertha Garlan]}|pwv} des Romans\pwindex{Schnitzler, Arthur 15.\,5.\,1862 Wien – 21.\,10.\,1931 ebd.@\textsc{Schnitzler, Arthur} (15.\,5.\,1862 Wien – 21.\,10.\,1931 ebd.), \emph{Schriftsteller, Mediziner}!Frau Bertha Garlan. Roman@\strich\emph{Frau Bertha Garlan. Roman}|pwv}, nachdem ich Sie von einem andern
                  Blatte\pwindex{?? [Zeitschrift, die Theodor von Sosnoskys Besprechung von »Bertha Garlan« zurückweist]@\emph{?? [Zeitschrift, die Theodor von Sosnoskys Besprechung von »Bertha Garlan« zurückweist]}|pw} als zu freundlich zurückerhalten, \label{K_L04230-2v}\edtext{an
               die »Nation\orgindex{Nation@Die Nation|pw}«}{\lemma{\textnormal{\emph{an
               die »Nation«}}}\Cendnote{\textnormal{Sosnoskys\pwindex{Sosnosky, Theodor von 4.\,1.\,1866 Budapest – 3.\,2.\,1943 Wien@\textsc{Sosnosky, Theodor von} (4.\,1.\,1866 Budapest – 3.\,2.\,1943 Wien), \emph{Schriftsteller}|pwk}{ }Rezension\pwindex{Sosnosky, Theodor von 4.\,1.\,1866 Budapest – 3.\,2.\,1943 Wien@\textsc{Sosnosky, Theodor von} (4.\,1.\,1866 Budapest – 3.\,2.\,1943 Wien), \emph{Schriftsteller}!Bücher und Zeitschriftenschau [Frau Bertha Garlan]@\strich\emph{Bücher und Zeitschriftenschau [Frau Bertha Garlan]}|pwkv} von \emph{Bertha Garland}\pwindex{Schnitzler, Arthur 15.\,5.\,1862 Wien – 21.\,10.\,1931 ebd.@\textsc{Schnitzler, Arthur} (15.\,5.\,1862 Wien – 21.\,10.\,1931 ebd.), \emph{Schriftsteller, Mediziner}!Frau Bertha Garlan. Roman@\strich\emph{Frau Bertha Garlan. Roman}|pwk} konnte erst viel später und an anderem Ort publiziert werden: Th. v. S.\pwindex{Sosnosky, Theodor von 4.\,1.\,1866 Budapest – 3.\,2.\,1943 Wien@\textsc{Sosnosky, Theodor von} (4.\,1.\,1866 Budapest – 3.\,2.\,1943 Wien), \emph{Schriftsteller}|pwk}: \emph{Bücher und Zeitschriftenschau}\pwindex{Sosnosky, Theodor von 4.\,1.\,1866 Budapest – 3.\,2.\,1943 Wien@\textsc{Sosnosky, Theodor von} (4.\,1.\,1866 Budapest – 3.\,2.\,1943 Wien), \emph{Schriftsteller}!Bücher und Zeitschriftenschau [Frau Bertha Garlan]@\strich\emph{Bücher und Zeitschriftenschau [Frau Bertha Garlan]}|pwk}. In: \emph{Norddeutsche Allgemeine Zeitung}\pwindex{Norddeutsche Allgemeine Zeitung@\emph{Norddeutsche Allgemeine Zeitung}|pwk}, Nr. 50,
                     28. 2. 1902, Beilage.}}}\label{K_L04230-2} geſandt; ſollte dieſe ſie nehmen,
              würde ich »Lt Gustel\pwindex{Schnitzler, Arthur 15.\,5.\,1862 Wien – 21.\,10.\,1931 ebd.@\textsc{Schnitzler, Arthur} (15.\,5.\,1862 Wien – 21.\,10.\,1931 ebd.), \emph{Schriftsteller, Mediziner}!Lieutenant Gustl. Novelle@\strich\emph{Lieutenant Gustl. Novelle}|pw}« dazunehmen. Dieſe Er{\pb}zählung\pwindex{Schnitzler, Arthur 15.\,5.\,1862 Wien – 21.\,10.\,1931 ebd.@\textsc{Schnitzler, Arthur} (15.\,5.\,1862 Wien – 21.\,10.\,1931 ebd.), \emph{Schriftsteller, Mediziner}!Lieutenant Gustl. Novelle@\strich\emph{Lieutenant Gustl. Novelle}|pwv}
               kenne ich bereits \label{K_L04230-3v}\edtext{aus der N. Fr. Pr\pwindex{Neue Freie Presse@\emph{Neue Freie Presse}|pw}}{\lemma{\textnormal{\emph{aus der N. Fr. Pr}}}\Cendnote{\textnormal{Arthur Schnitzler: \emph{Lieutenant Gustl}\pwindex{Schnitzler, Arthur 15.\,5.\,1862 Wien – 21.\,10.\,1931 ebd.@\textsc{Schnitzler, Arthur} (15.\,5.\,1862 Wien – 21.\,10.\,1931 ebd.), \emph{Schriftsteller, Mediziner}!Lieutenant Gustl. Novelle@\strich\emph{Lieutenant Gustl. Novelle}|pwk}. In: \emph{Neue Freie Presse}\pwindex{Neue Freie Presse@\emph{Neue Freie Presse}|pwk}, Nr. 13.053, 25. 12. 1900, Morgenblatt,
                     S. 34–41.}}}\label{K_L04230-3}. Sie wurde mir von befreundeter Seite zugeſchickt. Ich
               fand ſie hochintereßant, in der Technik brillant, will aber nicht verhehlen, daß sie
               in mir die heftigſte Opposition erweckt hat, da ich darin dieſelbe Feindſeligkeit
               gegen den Offiziersſtand erke{\geminationn}e, die mir aus »Freiwild\pwindex{Schnitzler, Arthur 15.\,5.\,1862 Wien – 21.\,10.\,1931 ebd.@\textsc{Schnitzler, Arthur} (15.\,5.\,1862 Wien – 21.\,10.\,1931 ebd.), \emph{Schriftsteller, Mediziner}!Freiwild. Schauspiel in 3 Akten@\strich\emph{Freiwild. Schauspiel in 3 Akten}|pw}« in Eri{\geminationn}erung ist. Wollen Sie \strikeout{mir} dieſe offene Ansicht, die
               kundzutun mir Pflicht ſcheint, nicht verübeln Ihrem Sie hochſchätzenden\pend
           \pstart \spacefill\mbox{ThvSosnosky}\pend{}\selectlanguage{ngerman}\endnumbering\briefempfaengerindex{Schnitzler, Arthur@\textsc{Schnitzler, Arthur}!zzzSosnosky, Theodor von@\emph{von Theodor von Sosnosky}!1901-05-231@{23. 5. 1901}|)be}\mylabel{L04230h}
\begin{anhang}
\end{anhang}\newcommand{\dateiname}{L04230}\newcommand{\titel}{Theodor von Sosnosky an Arthur Schnitzler, 23. 5. 1901}\newcommand{\editorInnen}{Herausgegeben von Jahnke, SelmaMüller, Martin Anton}\input{../tex-inputs/latex-leseansicht-abspann}
